\chapter*{Conclusion générale}
\addcontentsline{toc}{chapter}{Conclusion générale}
\markboth{Conclusion générale}{Conclusion générale}

Ce stage au sein du service Maintenance d'Air Liquide Tunisie a permis de concevoir et de réaliser AirLiquide GMAO+, une application qui remplace les saisies papier et tableur par un suivi unique de chaque appareil, de son retour à sa livraison. Les principaux apports sont une \textbf{fiche de vie unique} par appareil, une \textbf{traçabilité} horodatée et immuable, des \textbf{règles de workflow} appliquées par le serveur et vérifiées par 49 tests, ainsi qu'un journal de maintenance qui reprend le registre existant.

Sur le plan des compétences, le stage a mobilisé l'analyse de besoin avec des utilisateurs de terrain, la modélisation, le développement d'une API et d'une interface web, la sécurité applicative et la rédaction de tests. Il a aussi montré l'importance de comprendre le métier avant de coder : plusieurs choix, comme la séparation entre l'appareil et ses retours, viennent directement de l'observation du circuit réel.

Le travail présente des limites, détaillées au chapitre 4 : absence de rapports PDF, de maintenance préventive et de version mobile, et tests du frontend à compléter. Ces points constituent les perspectives naturelles du projet.

Sur le plan personnel, ce stage m'a permis de confronter ma formation de futur ingénieur en génie biomédical, parcours TTIS, à une situation de terrain : comprendre le circuit réel d'un dispositif médical dans un atelier de maintenance, puis le traduire en une application utilisable par les techniciens. J'y ai gagné en autonomie, en rigueur de documentation et en aisance pour dialoguer avec des utilisateurs dont le métier n'est pas l'informatique. Cette expérience confirme l'intérêt, pour un ingénieur biomédical, de maîtriser à la fois la maintenance des dispositifs médicaux et les outils numériques qui en assurent la traçabilité ; elle constitue une base solide pour mon projet professionnel dans la maintenance biomédicale et les systèmes d'information de santé.

